% gerado por ferramentas/md2tex_enunciado.py a partir de questoes/q091.md (não editar à mão)
O modelo de Rutherford para o átomo consiste em um núcleo bastante denso, pequeno e carregado positivamente, ao redor do qual orbitam partículas negativas. A interação atrativa entre o núcleo e os elétrons foi explicada posteriormente pela existência de uma força centrípeta de origem elétrica, cujo módulo, segundo a lei de Coulomb, é proporcional à constante eletrostática e ao quadrado da carga elementar, e inversamente proporcional ao quadrado da distância entre o núcleo e os elétrons. No entanto, o modelo planetário proposto por Rutherford não concorda com a teoria eletromagnética clássica, segundo a qual os elétrons descreveriam um movimento em espiral e se chocariam com o núcleo, o que não corresponde à realidade.\par
\figq[0.24\linewidth]{q091_fig1.png}
Qual gráfico apresenta a força de atração entre o núcleo e o elétron, em função do tempo para o átomo de hidrogênio, levando em consideração a teoria eletromagnética clássica?\par
\begin{center}
\altfig{A}{q091_altA.png}\hfill\altfig{B}{q091_altB.png}\hfill\altfig{C}{q091_altC.png}\hfill\altfig{D}{q091_altD.png}\hfill\altfig{E}{q091_altE.png}
\end{center}
